%==============================================================================
% Limitations + Future Work (~0.85 page)
% Written as a fair reviewer. Honest about single-org / small-N / single-seed /
% medium-once-audited / English-only / cost / gaming. Future work framed as OPEN
% research directions for the community, never as something the authors have
% built; no first-person "we are building", no naming any system.
%==============================================================================

\section{Limitations}
\label{sec:limitations}

We present this benchmark as a preliminary, single-seed release and are
deliberately conservative about what it establishes. The most important caveats
are the following.

\paragraph{A single synthetic organization.}
\label{para:single-synthetic-org}
Every result in this paper is computed over one fictional company, Helix
Logistics, generated by one open-weight model. The texture of a real
organizational corpus is controlled here for the sake of known ground truth
(Section~\ref{sec:design:generation}), but a single seed organization cannot
establish that the difficulty we observe generalizes to the messier statistics
of real companies---different industries, documentation cultures, team
topologies, and rates of decision churn. We also inherit the circularity
intrinsic to single-model generation: one model writes the artifacts, phrases
the questions, and runs the self-consistency check. The deterministic
graph-to-answer projection bounds label ambiguity but does not remove this risk,
and a cross-model validation pass remains future work. Following the datasheet
discipline of \citet{gebru2021datasheets}, we document these gaps rather than
obscure them throughout this section.

\paragraph{Small $N$, single seed, no human baseline yet.}
\label{para:small-n-single-seed}
The small tier has $11$ questions and the medium tier $73$, each scored once
with a single LLM judge, so we do not attach confidence intervals to individual
cells and do not claim that small between-system gaps are significant. Our own
replication probe is a warning. A host-versus-Docker re-run of an identically
configured system (\texttt{graphify-oss}) produced a $0.206$ vs.\ $0.220$
tier-mean pair---the \emph{average} essentially unchanged---while individual
\emph{per-question} scores swung by up to $\pm0.55$. The tier means are thus
stable enough to support the coarse ordering and the headline ``unsolved''
finding, but not precise enough to rank systems separated by a few points:
single-seed, per-question numbers are not defensible, and small-tier gaps below
roughly a tenth of a point should be read as indicative rather than significant
(the setup caveat is stated up front in \S\ref{sec:exp:setup}). We commit, for a
subsequent revision, to multi-seed runs with reported variance and bootstrap
confidence intervals, and to a human baseline on a question subset to anchor the
difficulty of the task independently of the automated judge.

\paragraph{Two tiers, once-audited labels.}
\label{para:two-tiers-manual-corrections}
We report two shipped tiers (small and medium), with all four public systems run
on both; the larger token tiers that would trace the scaling curve further are
not yet released. The medium tier was audited and repaired once against the
source corpus before the reported runs---the corrections change the labels, not
the systems, and are applied identically to every system---so a reader should
treat the medium numbers as resting on a once-audited label set rather than a
multiply-audited one. The specific audit-artifact classes, released pre-patch
snapshots, and per-question change logs are documented in
Appendix~\ref{app:annotation}. We caution, on the strength of the rank swap we
observe between tiers (Section~\ref{sec:exp:scaling}), against extrapolating
either tier's ordering to the unshipped larger tiers without measuring them.

\paragraph{English-only, one cost regime.}
\label{para:english-only-cost}
The corpus and questions are English-only; multilingual organizational memory,
where the same decision is discussed across languages, is unmeasured. The
evaluation itself is inexpensive per question, but cost compounds with the number
of seeds, systems, and tiers, and metered server-side extraction grows with
corpus size, so multi-seed sweeps across the larger tiers are the real budget
constraint and bias the scope of what others can replicate without a comparable
budget (the per-component and per-run dollar figures are reported in
Appendix~\ref{app:extended}).

\paragraph{Gaming risk.}
\label{para:gaming-risk}
Because the corpus is generated from a fixed taxonomy of six categories and a
known graph structure, a system tuned to that structure---for example, one that
hard-codes the bi-temporal field schema or the rubric facets---could inflate its
score without acquiring the underlying capability. The freshly generated seed
defends against memorization but not against structural overfitting. We mitigate
this with held-out tiers and paraphrased question phrasing, and we recommend that
leaderboard submissions report results on freshly generated organizations rather
than on the published seed alone.

\section{Future Work}
\label{sec:future}

\paragraph{}
\label{para:future-intro}
Beyond addressing the limitations above, our per-category failure analysis
points to concrete capabilities the field has not yet built. We frame these as
open research directions rather than a roadmap for any one system. The common
thread is that the categories where every public system collapses---bi-temporal
retrieval (C3), justification chains (C5), and contradiction (C6)---are exactly
the ones that flat, single-shot, document-granular memory cannot reach.

\paragraph{Richer entity-centric extraction.}
\label{para:future-entity-extraction}
The clearest architectural signal in our data is that coarse extraction breaks
at organizational scale: the weakest public system extracted roughly one graph
node per artifact (Appendix~\ref{app:extended}), so retrieval reached only
adjacent whole documents and the evidence never surfaced. A system that solved
C5 (justification chains) would likely need extraction at the granularity of
\emph{claims, actors, and events} rather than documents, so that the multi-hop
``why was this decided'' chain can be assembled from fine-grained nodes instead
of inferred from document adjacency.

\paragraph{Consolidation and background re-derivation.}
\label{para:future-consolidation}
As an organizational corpus grows, facts are superseded, corrected, and left in
unresolved contradiction faster than any single ingestion pass can reconcile. A
system that solved C1 (supersession) and C6 (contradiction) at scale would
plausibly need a background consolidation pass---a periodic ``dreaming'' step
that re-reads accumulated memory to re-derive which facts now supersede which,
attach explicit supersession reasons, and flag contradictions for surfacing
rather than silently resolving them---instead of relying solely on what was
inferable at write time.

\paragraph{Agentic, iterative retrieval.}
\label{para:future-agentic-retrieval}
The shift we observed at medium scale, from retrieval misses toward confident but
unsupported synthesis, suggests that single-shot retrieve-then-answer is the
wrong shape for these queries. A system that closed the gap on multi-hop,
cross-channel questions would likely need an iterative loop that retrieves,
inspects what it found, reformulates, and re-retrieves until the evidence chain
is complete or it can calibratedly abstain.

\paragraph{Explicit bi-temporal modeling.}
\label{para:future-bitemporal-modeling}
Every system in the field scored at or near the floor on the as-of-date
questions (C3, C4). A system that solved them would need a first-class
bi-temporal model---separating when a fact became true in the world from when
the organization recorded it, and supporting time-travel queries against a past
belief state---rather than treating ``as of'' as a hint to be ignored at
retrieval time. Designing benchmarks and systems around this distinction, which
existing conversational memory evaluations~\citep{maharana2024locomo,
wu2024longmemeval,beam2025} do not require, is in our view the central open
problem that organizational memory makes unavoidable.
